\chapter{Lecture 17: Exceptions and Design Patterns}

\tikzset{umlc/.style={draw,fill=blue!4,rectangle split,rectangle split parts=2,
    rectangle split part align={center,left},font=\small,inner sep=3pt,minimum width=2.6cm},
  inh/.style={-{Triangle[open,length=3mm,width=3mm]}}}
\begin{center}
\textit{March 8, 2018}
\end{center}

\section{Exceptions (continued)}
Last time: in C++, when an error condition occurs, an exception is thrown. By
default, an exception that is not handled causes the program to terminate.

To handle exceptions we use a \code{try \{ ... \}} block followed by
\code{catch (...) \{ ... \}} blocks (see \code{rangeErrorCaught.cc}):
\begin{lstlisting}[style=cpp]
#include <iostream>
#include <stdexcept>
#include <vector>
using namespace std;

int main() {
    vector<int> v;
    v.emplace_back(2);
    try {
        cout << v.at(1) << endl;    // out of range
    } catch (out_of_range &r) {
        cerr << "Range error: " << r.what() << endl;
    }
    cout << "Program continues" << endl;
}
\end{lstlisting}

\subsection{Stack Unwinding}
The thrower and the handler don't have to be in the same function
(\code{callChain.cc}):
\begin{lstlisting}[style=cpp]
void f() {
    cout << "Start f" << endl;
    throw out_of_range{"f"};        // the argument becomes what()
    cout << "Finish f" << endl;     // never printed
}
void g() { cout << "Start g" << endl; f(); cout << "Finish g" << endl; }
void h() { cout << "Start h" << endl; g(); cout << "Finish h" << endl; }

int main() {
    cout << "Start main" << endl;
    try {
        h();
    } catch (out_of_range &) {
        cerr << "Range error" << endl;
    }
    cout << "Finish main" << endl;
}
// Output: Start main, Start h, Start g, Start f, Range error, Finish main
\end{lstlisting}
\begin{center}
\begin{tikzpicture}[>=Stealth,font=\small,fr/.style={draw,minimum width=3.4cm,minimum height=0.6cm}]
  \node[fr,fill=blue!6] (m) at (0,0) {\code{main} (has \code{try}/\code{catch})};
  \node[fr,below=0pt of m,fill=gray!8] (h) {\code{h}};
  \node[fr,below=0pt of h,fill=gray!8] (g) {\code{g}};
  \node[fr,below=0pt of g,fill=red!10] (f) {\code{f}: \code{throw}};
  \node[font=\scriptsize,anchor=south] at (m.north) {call stack (grows down)};
  \draw[->,red!60!black,thick] (f.east) to[out=0,in=0,looseness=0.9] node[right,font=\scriptsize,align=left]{unwind: pop \code{f}, \code{g}, \code{h}\\ (running destructors of\\ their stack objects)} (m.east);
  \foreach \n in {h,g,f} \draw[gray] ([xshift=4pt]\n.west) -- ([xshift=-4pt]\n.east);
\end{tikzpicture}
\end{center}
``Finish f'', ``Finish g'' and ``Finish h'' are never printed: those frames are popped
without resuming.

\begin{definition}
\textbf{Stack unwinding}: when an exception is thrown, the call stack is
repeatedly popped until a matching \code{catch} block is found. Destructors of
stack-allocated objects in the popped frames run as their frames are popped.
\end{definition}

\subsection{Recovery in Stages}
A handler can do part of the recovery and then pass the problem on:
\begin{lstlisting}[style=cpp]
try {
    ...
} catch (SomeException &e) {
    // do part of the recovery, then choose:
    throw SomeOtherException{...};  // 1. throw a completely different exception
    throw e;                        // 2. throw the exception we caught
    throw;                          // 3. re-throw the original: recommended
}
\end{lstlisting}
In option 2, the exception we caught might really have been a \emph{subclass}
of \code{SomeException}; \code{throw e} throws a copy made from the
\emph{static} type of \code{e}, i.e.\ a sliced object. In option 3, the original
exception object is re-thrown as is; no slicing.

\subsection{What Can Be Thrown}
\begin{itemize}
  \item All exceptions in the C++ standard library inherit from the class
        \code{std::exception} (e.g.\ \code{out\_of\_range},
        \code{invalid\_argument}, \code{bad\_alloc}, \code{length\_error}).
  \item User-defined exceptions do not have to inherit from \code{exception}.
  \item In C++ you can throw anything: \code{int}s, strings, whatever.
\end{itemize}

Part of the standard exception hierarchy (all in \code{<stdexcept>} or \code{<new>}):
\begin{center}
\begin{tikzpicture}[font=\small,cls/.style={draw,fill=blue!4,minimum height=0.65cm,font=\ttfamily\small}]
  \node[cls] (e) at (0,0) {exception};
  \node[cls] (l) at (-3,-1.3) {logic\_error};
  \node[cls] (r) at (0.6,-1.3) {runtime\_error};
  \node[cls] (b) at (3.6,-1.3) {bad\_alloc};
  \node[cls] (o) at (-6,-2.6) {out\_of\_range};
  \node[cls] (i) at (-3.1,-2.6) {invalid\_argument};
  \node[cls] (le) at (-0.2,-2.6) {length\_error};
  \draw[inh] (l.north) -- (e.south west); \draw[inh] (r.north) -- (e.south); \draw[inh] (b.north) -- (e.south east);
  \draw[inh] (o.north) -- (l.south west); \draw[inh] (i.north) -- (l.south); \draw[inh] (le.north) -- (l.south east);
\end{tikzpicture}
\end{center}
A handler \code{catch (logic\_error \&e)} also catches \code{out\_of\_range}, since it is-a
\code{logic\_error}.

\paragraph{Q.} How do we write a catch block that catches all kinds of
exceptions?
\begin{lstlisting}[style=cpp]
try {
    ...
} catch (...) {     // literally three dots: catches everything
    ...
}
\end{lstlisting}

\subsection{Exceptions for Control Flow?}
Since exceptions jump out of arbitrarily deep calls, you could use them for
normal control flow, e.g.\ returning a result from deep recursion by throwing it.
\begin{quote}
\textit{``You can do some cool shit with exceptions, but don't use it.''}
--- Nomair Naeem
\end{quote}
Exceptions are slow to process: compare \code{recfib.cc} (ordinary recursive
Fibonacci) with \code{exfib.cc} (the same, but each call ``returns'' by
throwing): the throwing version is dramatically slower. Use exceptions for
actual errors.

\subsection{Good Practice}
Use the existing exception classes or create your own:
\begin{lstlisting}[style=cpp]
class BadInput {};

int n;
try {
    if (!(cin >> n)) throw BadInput{};
} catch (BadInput &) {              // naming the parameter (&b) is also legal
    cout << "Not an int, using 0" << endl;
    n = 0;
}
\end{lstlisting}
\begin{important}[Advice]
Always catch exceptions by reference:
\begin{itemize}
  \item no slicing occurs (subclass information, and virtual \code{what()},
        are preserved);
  \item it is more efficient, since no copy is made.
\end{itemize}
\end{important}

Exceptions to know: \code{vector::at} throws \code{std::out\_of\_range};
\code{new} throws \code{std::bad\_alloc} when there is no more heap memory.

\subsection{Exceptions and Destructors}
\begin{itemize}
  \item By default, if a destructor throws an exception, the program terminates
        (\code{std::terminate} is called), because dtors are implicitly
        \code{noexcept}.
  \item It is possible to write a destructor that may throw:
        \code{\textasciitilde Node() noexcept(false) \{ ... \}}. This is a bad
        idea.
\end{itemize}
Why: suppose the stack is unwinding because of one live exception. As frames are
popped, destructors run. If one of those destructors throws, there are now two
live exceptions, and C++ has no way to handle both: \code{std::terminate} is
called.

\section{Design Patterns (continued)}
The guiding principle behind design patterns: \textbf{program to an interface,
not an implementation}.
\begin{itemize}
  \item Use abstract base classes to provide the interface.
  \item Work through pointers/references to the abstract class and call its
        virtual methods.
  \item The actual objects can then be swapped in or out as needed: we get
        abstraction over a variety of behaviours.
\end{itemize}

\section{Iterator Pattern Revisited}
Our iterator implemented \code{*}, \code{++} and \code{!=}. Make that an
interface:
\begin{lstlisting}[style=cpp]
class AbsIter {
  public:
    virtual int &operator*() const = 0;
    virtual AbsIter &operator++() = 0;
    virtual bool operator!=(const AbsIter &other) const = 0;
    virtual ~AbsIter() {}
};

class List {
    struct Node;
    ...
  public:
    class Iterator : public AbsIter {
        // override the three pure virtual methods
    };
};

class Set {
    ...
  public:
    class Iterator : public AbsIter {
        ...
    };
};
\end{lstlisting}
Since \code{List::Iterator} and \code{Set::Iterator} both implement the
\code{AbsIter} interface, we can write generic traversals that work on either
data structure:
\begin{lstlisting}[style=cpp]
template <typename Fn>             // Fn: the type of a function parameter
void foreach(AbsIter &start, const AbsIter &finish, Fn f) {
    while (start != finish) {
        f(*start);
        ++start;
    }
}

void addFive(int &x) { x += 5; }

List l;
...
List::Iterator b = l.begin();
foreach(b, l.end(), addFive);      // adds 5 to every element of l
\end{lstlisting}
\begin{center}
\begin{tikzpicture}[font=\small,>=Stealth]
  \node[umlc,minimum width=4.2cm] (a) at (0,0) {\textit{AbsIter}\nodepart{second}\begin{tabular}{@{}l@{}}\code{+}\textit{operator*()}\\ \code{+}\textit{operator++()}\\ \code{+}\textit{operator!=()}\end{tabular}};
  \node[umlc,rectangle split parts=1] (l) at (-2.4,-2.4) {List::Iterator};
  \node[umlc,rectangle split parts=1] (s) at (2.4,-2.4) {Set::Iterator};
  \draw[inh] (l.north) -- (a.south west);
  \draw[inh] (s.north) -- (a.south east);
  \node[draw,rounded corners,fill=yellow!10] (fe) at (-5,0) {\code{foreach}};
  \draw[->,dashed] (fe) -- node[above,font=\scriptsize]{uses} (a);
\end{tikzpicture}
\end{center}

\section{The Observer Pattern}
Used for publish/subscribe systems:
\begin{itemize}
  \item the entity publishing data is the \textbf{subject} (publisher);
  \item the entities subscribing to the data are \textbf{observers}
        (subscribers).
\end{itemize}
The protocol:
\begin{enumerate}
  \item new data becomes available in the subject;
  \item the subject notifies all attached observers (\code{notifyObservers()}
        calls each observer's \code{notify()});
  \item each observer requests whatever further information it needs from the
        subject, e.g.\ by calling \code{getState()}.
\end{enumerate}
\begin{center}
\begin{tikzpicture}[>=Stealth,font=\small,o/.style={draw,rounded corners,fill=green!8,minimum width=2cm,minimum height=0.7cm}]
  \node[draw,rounded corners,fill=blue!6,minimum width=2.4cm,minimum height=1cm] (s) at (0,0) {subject};
  \node[o] (o1) at (5,0.9) {observer 1};
  \node[o] (o2) at (5,-0.9) {observer 2};
  \draw[->] ([yshift=3pt]s.east) -- node[above,sloped,font=\scriptsize]{2.\ \code{notify()}} ([yshift=3pt]o1.west);
  \draw[->] ([yshift=-3pt]s.east) -- node[below,sloped,font=\scriptsize]{2.\ \code{notify()}} ([yshift=-3pt]o2.west);
  \draw[->,dashed] (o1.north west) to[bend right=25] node[above,font=\scriptsize]{3.\ \code{getState()}} (s.north);
  \draw[->,dashed] (o2.south west) to[bend left=25] node[below,font=\scriptsize]{3.\ \code{getState()}} (s.south);
  \node[font=\scriptsize,anchor=east] at (s.west) {1.\ new data\ };
\end{tikzpicture}
\end{center}
Example: a spreadsheet's cells (subjects) and the charts and formulas that depend
on them (observers). Details and code next lecture.

\section*{Summary}
\begin{itemize}
  \item \code{try}/\code{catch}; stack unwinding pops frames (running dtors)
        until a handler is found.
  \item Re-throw with plain \code{throw;} to avoid slicing; catch by reference;
        \code{catch (...)} catches everything.
  \item Exceptions are for errors, not control flow (they are slow).
  \item Dtors should never throw: two simultaneous exceptions terminate the
        program.
  \item Program to interfaces; an abstract iterator enables generic algorithms.
  \item Observer: subjects notify attached observers of changes.
\end{itemize}
